\documentclass[../main]{subfiles}
\ifSubfilesClassLoaded{\setcounter{chapter}{5}}{}
\begin{document}

\chapter{Rigidbody, Input, dan Moving Player}

\begin{subcpmk}
  \item Sub-CPMK 3.2: Mengimplementasikan input player dan kontrol kamera
\end{subcpmk}

\section{Input System dan Player Movement}

\begin{learningoutcomes}
\item Menggunakan Unity Input System (Legacy/New)
\item Mengimplementasikan keyboard dan mouse input
\item Membuat player movement yang responsif
\end{learningoutcomes}

\subsection{Konsep Input di Unity}

Input Manager (Edit $\rightarrow$ Project Settings $\rightarrow$ Input Manager): Horizontal (A/D), Vertical (W/S), Fire1 (Ctrl, Mouse0), Jump (Space), Mouse X/Y.

\subsection{Implementasi Movement}

\begin{lstlisting}[language={[Sharp]C}]
void Update()
{
    float h = Input.GetAxis("Horizontal");
    float v = Input.GetAxis("Vertical");
    Vector3 move = new Vector3(h, v, 0) * speed * Time.deltaTime;
    transform.Translate(move);
}
\end{lstlisting}

\subsection{Camera Follow}

Gunakan LateUpdate() untuk follow player. Atur posisi Main Camera mengikuti target.

\section{Shooting dan Camera Control}

\begin{learningoutcomes}
\item Mengimplementasikan shooting mechanic
\item Mengatur camera follow dan control
\item Menggunakan touch/gamepad untuk mobile
\end{learningoutcomes}

\subsection{Shooting}

Input.GetButtonDown("Fire1") untuk satu tembakan. Instantiate(bulletPrefab) untuk spawn bullet. Atur posisi dan velocity bullet.

\subsection{Camera Control}

Cinemachine package untuk camera follow otomatis. Atau manual: kamera posisi = target posisi + offset.

\subsection{Platform Input}

Touch input: Input.touchCount, Input.GetTouch(). Gamepad: Input.GetAxis("Horizontal") dengan joystick. New Input System untuk fleksibilitas lebih.


\begin{aktivitas}
  \item Pelajari konsep Rigidbody, AddForce, dan Input.GetAxis melalui materi di atas
  \item Diskusikan mengapa AddForce harus dipanggil di FixedUpdate
  \item Diskusi konsep AddForce dan FixedUpdate
\end{aktivitas}

\begin{latihan}
  \item Mengapa AddForce dipanggil di FixedUpdate, bukan Update?
  \item Apa peran Input.GetAxis dalam kontrol player?
\end{latihan}

\begin{checklist}
  \item Saya dapat menjelaskan konsep Rigidbody dan AddForce
  \item Saya memahami penggunaan FixedUpdate untuk fisika
\end{checklist}

\begin{rangkuman}
Bab ini membahas teori gerakan player dengan Rigidbody, \method{AddForce}, dan \method{Input.GetAxis}. Contoh kode sebagai ilustrasi konsep.
\end{rangkuman}

\ifSubfilesClassLoaded{
  \renewcommand{\bibname}{Daftar Pustaka}
  \bibliographystyle{plain}
  \bibliography{../references}
}{}
\end{document}
